\chapter{Intermezzo: Connes' noncommutative geometry}\label{ch:ncg}

Geometry from algebra: spectral triples, the Dixmier trace, the spectral action principle, almost-commutative geometries and the Standard Model, the Lorentzian case, thermal time, and a speculative bridge to the rest of the book.

\section{Geometry from algebra}\label{ch:gelfand}
\skippable

\begin{physmotiv}
A classical phase space is a set of points on which observables are functions; the functions form a commutative algebra. Quantization replaces this algebra by a noncommutative one, and in that sense phase space itself becomes a ``noncommutative space''. Connes' noncommutative geometry takes this literally: a space \emph{is} its algebra of functions. This section explains why that is a theorem in the commutative case and what it suggests in general. Part IX is a self-contained detour; it connects back to Part IV through the modular flow (\cref{ch:thermal_time}).
\end{physmotiv}

\subsection{Gelfand--Naimark duality}

\begin{theorem}[Gelfand--Naimark]
Every commutative unital \cstar-algebra $\A$ is isomorphic to $C(X)$ for a compact Hausdorff space $X$, namely the space of characters (nonzero $^*$-homomorphisms $\A\to\C$, equivalently pure states) with the weak-$^*$ topology. Continuous maps $X\to Y$ correspond to $^*$-homomorphisms $C(Y)\to C(X)$ (contravariantly).
\end{theorem}

So compact spaces and commutative \cstar-algebras are the same thing, with points $=$ pure states (\cref{ch:states}: pure states of $C(X)$ are Dirac measures). Similarly, measure spaces correspond to commutative von Neumann algebras $\Linf(X,\mu)$, and the dictionary extends: \emph{noncommutative \cstar-algebras are noncommutative topological spaces; noncommutative von Neumann algebras are noncommutative measure spaces}.

\subsection{What is gained}
\begin{itemize}
\item \textbf{Singular quotients become algebras.} The quotient of a space by a group action that is not free (e.g.\ irrational rotation of the circle) is topologically pathological, but the crossed product $C(S^1)\rtimes\Z$ is a perfectly good \cstar-algebra, the ``noncommutative torus''.
\item \textbf{Phase space.} The Weyl algebra of \cref{ch:weyl} is the noncommutative torus algebra when $x,p$ are periodic; $\hbar$ is the noncommutativity parameter.
\item \textbf{Crossed products are spaces too.} The crossed product $\M\rtimes\R$ of \cref{ch:crossed} is the noncommutative version of the quotient of a ``space'' by an $\R$-action: for a measurable space with an action it is the groupoid algebra of the orbit equivalence relation. So the passage III$\to$II$_\infty$ is a geometric operation: ``forming the transversal''.
\end{itemize}

\begin{keyidea}
Commutative algebras are spaces. Noncommutative algebras are generalized spaces. Constructions on spaces (quotients, product, bundles) have algebraic counterparts (crossed products, tensor products, modules).
\end{keyidea}

\subsection*{Exercises}
\begin{exercise}\label{ex:gelfand:1}
Show the characters of $C(X)$ are evaluations $f\mapsto f(x)$. (Hint: the kernel of a character is a maximal ideal; if no point $x$ had $f(x)=0$ for all $f$ in the ideal, compactness and a partition of unity would put an invertible element in it.)
\end{exercise}
\begin{exercise}\label{ex:gelfand:2}
Show that $C(S^1)\rtimes_\theta\Z$ for the rotation by an angle $2\pi\theta$ is generated by unitaries $U,V$ with $UV=\ee^{2\pi\ii\theta}VU$ (noncommutative torus) and compare with the Weyl relations (\cref{ch:weyl}).
\end{exercise}

\section{Spectral triples}\label{ch:spectral_triples}
\skippable

\begin{physmotiv}
Topology is the algebra of continuous functions; \emph{metric geometry} needs one more ingredient. Connes' observation is that the Dirac operator supplies it: distances are encoded in how fast functions can vary, measured by commutators $[D,f]$. A spectral triple $(\A,\HH,D)$ is a noncommutative Riemannian (spin) manifold.
\end{physmotiv}

\subsection{Definition}

\begin{definition}
A (unital, even/odd) \emph{spectral triple}\index{spectral triple} $(\A,\HH,D)$ consists of a $^*$-algebra $\A$ represented on a Hilbert space $\HH$ and a self-adjoint operator $D$ on $\HH$ with compact resolvent such that $[D,a]$ is bounded for all $a\in\A$. (Further axioms: $\Z_2$-grading, real structure $J$ (with sign conventions depending on a KO-dimension), orientability, first order condition, regularity, Poincar\'e duality.)
\end{definition}

\textbf{Canonical example.} $\A=C^\infty(M)$, $\HH=L^2(M,S)$ spinors on a compact Riemannian spin manifold $M$, $D=\ii\gamma^\mu\nabla_\mu$ the Dirac operator. Then $[D,f]=\ii\gamma^\mu\partial_\mu f$ and $\norm{[D,f]}=\norm{\nabla f}_\infty$ (Lipschitz seminorm).

\subsection{Distance formula and reconstruction}

The geodesic distance on $M$ is recovered from the triple:
\begin{equation}
d(p,q)=\sup\big\{|f(p)-f(q)|:f\in\A,\ \norm{[D,f]}\le1\big\}.
\end{equation}
This makes sense for any triple, and defines a distance on the \emph{state space} of $\A$ (Connes' metric), the noncommutative substitute for points.

\begin{theorem}[Connes' reconstruction theorem \cite{Connes2013}]
A commutative spectral triple satisfying Connes' axioms comes from a compact Riemannian spin$^c$ manifold. (Connes 2008, 2013.)
\end{theorem}

\emph{Physics hint.} The distance formula is the statement that distance is the maximum change of a ``potential'' with bounded gradient; $D$ is the algebraic gradient, as the Dirac operator is the ``square root'' of the Laplacian.

\begin{whatwrong}
Spectral triples are \emph{Euclidean}: $D$ is self-adjoint with a spectrum, and the distance formula needs positivity. Lorentzian generalizations are an open area (\cref{ch:lorentzian_nc}).
\end{whatwrong}

\subsection*{Exercises}
\begin{exercise}\label{ex:spectral_triples:1}
For $M=S^1$ of length $L$ and $D=-\ii\,\dd/\dd x$, compute $d(p,q)$ from the formula and recover the geodesic distance. (Take $f(x)$ with $|f'|\le1$.)
\end{exercise}
\begin{exercise}\label{ex:spectral_triples:2}
Two-point space: $\A=\C^2$, $\HH=\C^2$, $D=\begin{pmatrix}0&m\\m&0\end{pmatrix}$ (with $\gamma$-grading). Compute $[D,f]$ for $f=(f_1,f_2)$ and show $d(1,2)=1/m$.
\end{exercise}

\section{Infinitesimals and the Dixmier trace}\label{ch:dixmier}
\skippable

\begin{physmotiv}
Integration in noncommutative geometry has to be defined algebraically. Connes' idea: compact operators are \emph{infinitesimals}, the order of an infinitesimal is the decay rate of its eigenvalues, and the \emph{Dixmier trace}\index{Dixmier trace} extracts the coefficient of the logarithmic divergence of the partial traces, which in the manifold case is the Riemannian volume integral. It is a trace on a \emph{ideal} of compact operators; it is a singular trace on $\Bnd{\HH}$ and so related to the traces on type II factors of \cref{ch:traces}.
\end{physmotiv}

\subsection{Infinitesimals}

A compact operator $T$ has singular values $\mu_n(T)\downarrow0$. We say $T$ is an infinitesimal of order $\alpha$ if $\mu_n(T)=O(n^{-\alpha})$. For a $d$-dimensional manifold, the Weyl law implies that $|D|^{-1}$ has $\mu_n\sim n^{-1/d}$, so $|D|^{-1}$ is an infinitesimal of order $1/d$: the ``line element $\dd s$'', and $\dd s^d=|D|^{-d}$ is of order $1$, the borderline of trace class.

\subsection{The Dixmier trace}

For $T\ge0$ in the ideal $\mathcal L^{1,\infty}$ (partial sums $\sum_{n\le N}\mu_n=O(\log N)$) set
\begin{equation}
\Tr_\omega(T)=\mathrm{Lim}_\omega\frac1{\log N}\sum_{n=1}^N\mu_n(T),
\end{equation}
where $\mathrm{Lim}_\omega$ is a generalized limit (Hahn--Banach). $\Tr_\omega$ is linear, positive, vanishes on trace-class operators and is a \emph{trace}: $\Tr_\omega(ST)=\Tr_\omega(TS)$, and it is independent of $\omega$ on measurable operators.

\begin{theorem}[Connes' trace theorem]
On a compact $d$-manifold, for a pseudodifferential operator $P$ of order $-d$, $\Tr_\omega(P)=\frac1d\,\mathrm{Res}_W(P)$ (the Wodzicki residue). For $P=f|D|^{-d}$: $\Tr_\omega\big(f|D|^{-d}\big)=c_d\int_Mf\,\dd\mathrm{vol}$ with $c_d$ a universal constant.
\end{theorem}

Thus \emph{the Riemannian volume integral is a Dixmier trace}: $\int f=\Tr_\omega(f|D|^{-d})$ up to a constant. This is the noncommutative integral.

\subsection{Relation to traces on type II factors}

A trace on a type II$_1$ factor with values in $[0,1]$ on projections measures ``dimension'' continuously (\cref{ch:classification}). The Dixmier trace is a different thing (singular trace on compact operators), but they are related in spirit: both are positive traces extracting a regularized ``size''. In noncommutative geometry applied to foliations, the transverse measure is a type II trace on the von Neumann algebra of the foliation, whose Dixmier-type integration gives transverse volume; the interplay with the Tomita--Takesaki modular flow appears for type III (non-invariant measures).

\subsection*{Exercises}
\begin{exercise}\label{ex:dixmier:1}
For $T=\mathrm{diag}(1,\tfrac12,\tfrac13,\dots)$ show $\sum_{n\le N}\mu_n\sim\log N$, so $T\in\mathcal L^{1,\infty}$ with $\Tr_\omega T=1$.
\end{exercise}
\begin{exercise}\label{ex:dixmier:2}
For $S^1$ with $D=-\ii\,\dd/\dd x$ (eigenvalues $n\in\Z$ for length $2\pi$), compute $\Tr_\omega(|D|^{-1})$ on the non-zero modes and compare with $\int_{S^1}\dd x/2\pi\times$ const.
\end{exercise}

\section{The spectral action principle}\label{ch:spectral_action}
\skippable

\begin{physmotiv}
What action should one use in a world where the only data are the spectral triple $(\A,\HH,D)$? Chamseddine and Connes \cite{CCM1997} proposed the simplest invariant: $S=\Tr f(D/\Lambda)$, a function of the spectrum of $D$ alone, with $\Lambda$ a cutoff scale. Its asymptotic expansion in $\Lambda$ reproduces the cosmological constant, the Einstein--Hilbert term and higher-curvature terms. ``One can hear the shape of the universe'': the action depends only on the spectrum.
\end{physmotiv}

\subsection{The principle and the heat kernel expansion}

\begin{equation}
S[D]=\Tr f\!\Big(\frac{D}{\Lambda}\Big)+\text{(fermionic action }\inner{\psi}{D\psi}),
\end{equation}
with $f$ a positive even cutoff function. It is diffeomorphism invariant (depends only on the unitary class of $D$). For a four-dimensional compact Riemannian spin manifold, the large-$\Lambda$ asymptotics follows from the Seeley--DeWitt heat kernel expansion of $\Tr\ee^{-tD^2}\sim\sum t^{(k-4)/2}a_k(D^2)$:
\begin{equation}
\Tr f\!\Big(\frac D\Lambda\Big)\sim f_4\Lambda^4a_0+f_2\Lambda^2a_2+f_0a_4+\cdots,
\end{equation}
where $f_k$ are moments of $f$ and, for $D$ the Dirac operator,
\begin{align}
a_0&\propto\int\dd^4x\sqrt g\quad\text{(cosmological constant)},\\
a_2&\propto\int\dd^4x\sqrt g\,R\quad\text{(Einstein--Hilbert, with $G_N\sim1/f_2\Lambda^2$)},\\
a_4&\propto\int\dd^4x\sqrt g\big(c_1C^2_{\mu\nu\rho\sigma}+c_2\,\chi\big)\quad\text{(Weyl}^2\text{ and Euler/Gauss--Bonnet)}.
\end{align}
(Coefficients depend on conventions; the structure is robust.) So the spectral action contains gravity, with the three parameters $f_4\Lambda^4,f_2\Lambda^2,f_0$ determined by the moments of one function.

\begin{whatwrong}
This is a Euclidean action on a compact manifold; its relation to the Lorentzian physical action requires Wick rotation, problematic in a general background (\cref{ch:lorentzian_nc}). The values of the couplings are those at the cutoff scale $\Lambda$, and need RG running to the physical scales; this is the source of the Standard Model predictions and their problems (\cref{ch:almost_comm}).
\end{whatwrong}

\subsection*{Exercises}
\begin{exercise}\label{ex:spectral_action:1}
For the flat torus $T^4$ compute $\Tr\ee^{-tD^2}$ for the Dirac operator and show $a_0\propto\mathrm{vol}$ and all higher $a_k$ vanish.
\end{exercise}
\begin{exercise}\label{ex:spectral_action:2}
Show that for $f(x)=\ee^{-x^2}$, $f_4,f_2,f_0$ are proportional to $\int x^3f$, $\int xf$ and $f(0)$; compute them.
\end{exercise}

\section{Almost-commutative geometries and the Standard Model}\label{ch:almost_comm}
\skippable

\begin{physmotiv}
Take spacetime times a tiny \emph{finite} noncommutative space. The Dirac operator acquires components along the finite part; gauge fields come from the manifold directions and the Higgs field from the discrete directions, as parts of a single ``connection''. With the right finite algebra, the gauge group and fermion representations of the Standard Model come out. This is the Chamseddine--Connes--Marcolli picture. It is elegant, makes predictions, and has an honest, mixed track record, which we describe.
\end{physmotiv}

\begin{unsettled}
The Standard Model reconstruction is a research program with contested status. Statements about predictions and their fate below are summarized from memory of the literature and must be verified (see \texttt{refs.bib}).
\end{unsettled}

\subsection{Product geometries}

An \emph{almost-commutative} geometry is the product of a manifold triple $(C^\infty(M),L^2(M,S),D_M)$ and a finite triple $(\A_F,\HH_F,D_F)$: algebra $C^\infty(M)\otimes\A_F$, Hilbert space $L^2(M,S)\otimes\HH_F$, $D=D_M\otimes\id+\gamma_5\otimes D_F$. The finite algebra $\A_F$ is a sum of matrix algebras. \emph{Inner fluctuations} $D\mapsto D+A+JAJ^{-1}$ with $A=\sum a_i[D,b_i]$ produce: gauge bosons from the $M$-components of $A$, and scalar fields (Higgs) from the $F$-components, of the unitary group of the algebra; $D_F$ encodes Yukawa couplings and masses.

\subsection{The finite algebra of the Standard Model}

The choice $\A_F=\C\oplus\mathbb H\oplus M_3(\C)$ (with $\mathbb H$ the quaternions) acting on a space whose basis are the fundamental fermions and their antiparticles reproduces the $SU(3)\times SU(2)\times U(1)$ structure and the hypercharges after imposing the first-order condition and a KO-dimension $6\bmod8$ finite geometry (Chamseddine--Connes). The spectral action then yields the bosonic Lagrangian: Yang--Mills terms, Higgs potential, and gravity, with relations among the couplings at the cutoff scale.

\subsection{What is predicted, and what is not}

\begin{itemize}
\item \textbf{Relations at the unification scale:} gauge coupling unification relations (like $g_2^2=g_3^2=\tfrac53g_1^2$) and a relation between the Higgs quartic, top Yukawa and gauge couplings. Running these down to low energies predicted a Higgs mass of order 170 GeV (Chamseddine--Connes--Marcolli). The measured 125 GeV is in tension with this prediction; proposed repairs include an additional real scalar field coupled to the Higgs and neutrino sector (Chamseddine--Connes--van Suijlekom \cite{CCvS2013} and others) and other modifications (e.g.\ departures from the original assumptions about the Dirac operator). \textbf{Status: contested.}
\item \textbf{Number of generations and gravity:} the finite space accommodates three generations as input, not a derivation.
\item \textbf{Neutrinos:} Majorana-type terms and the see-saw can be incorporated.
\end{itemize}

\begin{whatwrong}
This is not a proof that the Standard Model ``is'' noncommutative geometry; a different finite geometry, or the failure of the unification boundary conditions, changes the predictions. The program is best seen as a constraint-satisfaction exercise: a small set of axioms (spectral triple axioms) selects a restricted family of field content and relations.
\end{whatwrong}

\subsection*{Exercises}
\begin{exercise}\label{ex:almost_comm:1}
For $\A_F=\C\oplus\C$, $\HH_F=\C^2$, $D_F=\begin{pmatrix}0&m\\m&0\end{pmatrix}$ compute the inner fluctuations $A=\sum a[D_F,b]$ and show they form a single complex scalar field $\phi$ with $D_F\to\begin{pmatrix}0&m+\phi\\\bar m+\bar\phi&0\end{pmatrix}$; compute the spectral action $\Tr f(D_F/\Lambda)$ and identify a quartic potential for $\phi$ (a toy Higgs).
\end{exercise}
\begin{exercise}\label{ex:almost_comm:2}
Explain how the Lie algebra of the unitary group of $\C\oplus\mathbb H\oplus M_3(\C)$ contains $\mathfrak u(1)\oplus\mathfrak{su}(2)\oplus\mathfrak u(3)$ and how unimodularity cuts $\mathfrak u(3)\to\mathfrak{su}(3)$ and constrains hypercharges.
\end{exercise}

\section{Lorentzian signature and open problems}\label{ch:lorentzian_nc}
\skippable

\begin{physmotiv}
Physics is Lorentzian; spectral triples are Euclidean. This section lists what is known and not about moving to Lorentzian signature, and other open issues of the program.
\end{physmotiv}

\begin{itemize}
\item \textbf{Krein spaces.} A Lorentzian Dirac operator is not self-adjoint on a Hilbert space but is self-adjoint with respect to an indefinite inner product on a Krein space; several Lorentzian spectral triple proposals exist, with axioms differing from the Euclidean ones.
\item \textbf{Wick rotation.} In general no analytic continuation exists on a time-dependent background; the spectral action in its heat-kernel form is Euclidean. A spectral \emph{Lorentzian} action and its asymptotic expansion are not established in generality.
\item \textbf{Causal structure.} Distance formulas have Lorentzian analogues (Moretti, Franco--Eckstein, and others) involving a causal cone of states; no consensus.
\item \textbf{Quantization.} The spectral action is classical; coupling to quantum matter fields and gravitational quantization remains unclear.
\item \textbf{Unitarity and positivity} of the resulting field theory with higher-derivative terms from $a_4$.
\end{itemize}
(These remarks summarize the field's open directions and should be checked against surveys.)

\subsection*{Exercises}
\begin{exercise}\label{ex:lorentzian_nc:1}
Why does the heat kernel $\Tr\ee^{-tD^2}$ fail to exist for a Lorentzian Dirac operator? (Hint: $D^2$ is not bounded below.)
\end{exercise}

\section{Thermal time (Connes--Rovelli)}\label{ch:thermal_time}
\skippable

\begin{physmotiv}
In a generally covariant theory there is no preferred time. Connes and Rovelli proposed that time is \emph{emergent from the state}: for a state $\omega$ of the observable algebra, the \emph{thermal time} is the modular flow $\sigma^\omega_t$. The Connes cocycle theorem makes this canonical up to inner automorphisms, so the \emph{outer} class of the flow is a state-independent notion of ``time'', a hidden time inherent to the algebra. Here we contrast it with the observer's clock of \cref{ch:clocks_qrf}.
\end{physmotiv}

\subsection{The hypothesis}

\textbf{Thermal time hypothesis} (Connes--Rovelli 1994 \cite{ConnesRovelli1994}): the physical time flow is the modular flow $\sigma^\omega_t$ of the state $\omega$ in which the system is found. By \cref{ch:modham} different faithful states give flows differing by inner automorphisms, so a canonical outer flow $\delta:\R\to\mathrm{Out}(\M)$ exists. For a Gibbs state with Hamiltonian $H$: $\sigma_t=\Ad\ee^{-\ii\beta Ht}$: thermal time is $-\beta\times$ ordinary time (temperature as the ratio of time units). For the Rindler vacuum: the boost with the Unruh temperature: Unruh temperature as the ratio of proper time to modular time.

\subsection{Comparison with the observer's clock}

\begin{center}
\renewcommand{\arraystretch}{1.3}
\resizebox{\ifdim\width>\linewidth\linewidth\else\width\fi}{!}{%
\begin{tabular}{@{}p{3.2cm}p{4.8cm}p{4.8cm}@{}}
\toprule
 & Thermal time & Observer's clock\\
\midrule
Source & the state $\omega$ & a physical subsystem with Hamiltonian $q$\\
Flow & modular, $\sigma_t^\omega$ & dressing by $\ee^{-\ii pK}$ w.r.t.\ clock reading\\
Canonical? & outer class only & depends on the clock chosen\\
Role in this book & generates the crossed product & provides the extra variable $x$\\
\bottomrule
\end{tabular}}
\end{center}
The crossed product combines the two: the modular flow of the state provides the automorphism, the observer's clock provides the variable that makes it inner (\cref{ch:add_observer}). The thermal time hypothesis is thus ``half'' of the CLPW construction: it says what the flow \emph{is}; the observer says whom it is time \emph{for}.

\begin{whatwrong}
The hypothesis is conceptual and has been criticized (a modular flow is determined by a state, which seems to presuppose the dynamics it explains; and different states give genuinely different flows up to inner automorphisms). Its precise status is not settled.
\end{whatwrong}

\subsection*{Exercises}
\begin{exercise}\label{ex:thermal_time:1}
For the thermofield double, identify thermal time with $-\beta t$ and compute the relation between Unruh and Hawking temperatures as modular-to-proper time ratios.
\end{exercise}
\begin{exercise}\label{ex:thermal_time:2}
Show that for a trace, the modular flow is trivial, so thermal time is frozen; interpret as the statement that infinite temperature has no time.
\end{exercise}

\section{Speculative bridge}\label{ch:bridge}
\skippable

\begin{physmotiv}
The algebraic-gravity program (Parts V--VII) and spectral geometry (Part IX) both treat algebras as fundamental and both rely on traces and modular theory. This section lists possible bridges. It is \emph{speculative} and should be read as questions, not results.
\end{physmotiv}

\begin{unsettled}
Everything in this section is speculation or an open direction, unless explicitly attributed.
\end{unsettled}

\begin{enumerate}
\item \textbf{Traces.} Both programs integrate with a trace: the type II$_1$ trace of \cref{ch:clpw} (entropy) and the Dixmier trace of \cref{ch:dixmier} (volume). Could the entropy of a horizon be expressed as a noncommutative integral of a spectral triple associated with the static patch?
\item \textbf{Time.} Thermal time (\cref{ch:thermal_time}) and the observer's clock (\cref{ch:clocks_qrf}) are two notions of time from algebras; a unified treatment (the spectral triple of the observer algebra) is not known.
\item \textbf{Spectral data of horizons.} Is there a canonical spectral triple for a Killing horizon whose Dirac operator is the modular Hamiltonian?
\item \textbf{Almost-commutative structure of the observer.} The observer's internal algebra $\A_{\rm obs}$ (finite) combined with spacetime: product geometry or crossed product?
\end{enumerate}

\subsection*{A bonus: the Connes embedding problem}
Connes' embedding conjecture asked whether every separable type II$_1$ factor embeds in an ultrapower of the hyperfinite II$_1$ factor $R$, i.e.\ whether the hyperfinite factor is a universal approximant. It was shown to be \emph{false} as a consequence of the result $\mathrm{MIP}^*=\mathrm{RE}$ of Ji, Natarajan, Vidick, Wright and Yuen \cite{MIPstarRE} (2020), connecting quantum information (non-local games, Tsirelson's problem) with operator algebras. Relevance here: the type II$_1$ algebra $\A_{\rm dS}$ of \cref{ch:clpw}, if the static patch were described by a finite-dimensional approximant, would be embeddable in $R^\omega$; whether the physical $\A_{\rm dS}$ is embeddable is an interesting question that now has no a priori answer.

\subsection*{Exercises}
\begin{exercise}\label{ex:bridge:1}
State precisely what ``$\A$ embeds in $R^\omega$'' means via matrix approximations $\tr(a_{i_1}\cdots a_{i_k})\approx\frac1{n}\Tr(A_{i_1}\cdots A_{i_k})$ for finite matrices $A_i$.
\end{exercise}


\section{Two computations: the Dirac spectrum on the sphere and the fermion count}\label{ch:ncgcomp}
\begin{physmotiv}
Two checks make the spectral-action machinery concrete. A spectrum that can be summed by hand confirms the first two heat-kernel coefficients, and a head count of the fermions confirms that the finite triple of the Standard Model has the right size.
\end{physmotiv}

\subsection{The Dirac operator on $S^2$}
On the unit two-sphere the Dirac operator has eigenvalues $\pm m$, $m=1,2,3,\dots$, each with multiplicity $2m$. So
\begin{equation}
\Tr e^{-tD^2}=\sum_{m\ge1}4m\,\ee^{-tm^2}.
\end{equation}
Euler--Maclaurin for $f(m)=m\,\ee^{-tm^2}$ ($f(0)=0$, $f'(0)=1$) gives $\sum_{m\ge1}f(m)=\frac1{2t}-\frac1{12}+O(t)$, hence
\begin{equation}
\Tr e^{-tD^2}=\frac2t-\frac13+O(t).
\end{equation}
Compare with the heat-kernel expansion $\Tr e^{-tD^2}\simeq\frac1{4\pi t}\int\mathrm{tr}\big(1+t(\tfrac R6-\tfrac R4)+\dots\big)$ for the Dirac operator squared, $D^2=\nabla^*\nabla+R/4$: with $\mathrm{tr}\,1=2$, $R=2$ and area $4\pi$, the leading term is $2/t$ and the next is $-\frac1{12}R\cdot2\cdot\frac{4\pi}{4\pi}=-\frac13$. They agree. This is the same $a_0$ and $a_2$ that appear in the spectral action (\cref{ch:spectral_action}) as the cosmological and Einstein--Hilbert terms.

\subsection{Counting fermions}
In one generation of the Standard Model with a right-handed neutrino the Weyl fermions are $\nu_L,e_L,\nu_R,e_R$ (4 leptons) and $u_L,d_L,u_R,d_R$ in three colours (12 quarks): 16 in total. Adding the antiparticles doubles it and three generations triple it:
\begin{equation}
\dim\HH_F=3\times2\times16=96.
\end{equation}

\subsection*{Exercises}
\begin{exercise}\label{ex:ncgcomp:1}
Check the multiplicities $2m$ for the lowest level $m=1$ by counting the Killing spinors on $S^2$, or derive the whole spectrum from $D^2=-\Delta_{S^2}+\ldots$ on the spinor bundle.
\end{exercise}
\begin{exercise}\label{ex:ncgcomp:2}
Carry the Euler--Maclaurin expansion one order further to get the $O(t)$ coefficient of $\Tr e^{-tD^2}$, and check your answer numerically at $t=0.05$.
\end{exercise}
\begin{exercise}\label{ex:ncgcomp:3}
Show that the unitary group of the quaternions $\mathbb H$ (norm-one elements) is $SU(2)$ and that $U(M_3(\C))=U(3)$. Explain how the Standard Model gauge group arises from $\C\oplus\mathbb H\oplus M_3(\C)$ after the unimodularity condition.
\end{exercise}
\begin{exercise}\label{ex:ncgcomp:4}
Recount the 96 degrees of freedom without a right-handed neutrino and find $\dim\HH_F$. Which part of the finite Dirac operator is lost?
\end{exercise}
